\section{Трассировка PDOM\_Diverge на примере вложенной дивергенции}
\label{app:pdom_trace}

Рассмотрим фрагмент CUDA-кода с двумя вложенными условными ветвлениями:

\begin{lstlisting}[language=C++, caption={Пример вложенной дивергенции}]
// tid = threadIdx.x
if (tid < 16) {           // outer if: threads 0-15 / 16-31
    if (tid < 8) {        // inner if: threads 0-7 / 8-15
        A();
    } else {
        B();
    }
    C();
} else {
    D();
}
E();  // reconvergence point
\end{lstlisting}

Начальная маска: $M_0 = \texttt{0xFFFFFFFF}$.
Базовые блоки: $\mathrm{bb\_outer\_if}$, $\mathrm{bb\_inner\_if}$,
$\mathrm{bb\_A}$, $\mathrm{bb\_B}$, $\mathrm{bb\_C}$, $\mathrm{bb\_D}$,
$\mathrm{bb\_E}$ (реконвергенция).

Пусть $\mathrm{PC}_E = 100$~--- PC метки \texttt{E},
$\mathrm{PC}_C = 50$~--- PC метки \texttt{C}.

\begin{table}[h]
  \caption{Трассировка стека дивергенции}
  \label{tab:pdom_trace}
  \centering
  \begin{tabular}{clp{8cm}}
    \toprule
    \textbf{Шаг} & \textbf{Текущий PC / инструкция} & \textbf{Стек $S$ (вершина справа)} \\
    \midrule
    1 & \texttt{setp.lt} + \texttt{@\%p bra \$outer}
      & Внешняя дивергенция: \\
    & & $M_{\mathrm{taken}} = \texttt{0x0000FFFF}$,
        $M_{\mathrm{not}} = \texttt{0xFFFF0000}$ \\
    & & push конвергенции $\{100, \texttt{0xFFFFFFFF}\}$,
        push «не взята» $\{PC_{D}, \texttt{0xFFFF0000}\}$ \\
    & & $S = [\{100, M_0, T\},\; \{PC_D, M_{\mathrm{not}}, F\}]$ \\
    \midrule
    2 & Исполнение $\mathrm{bb\_inner\_if}$: \texttt{@\%p bra \$inner}
      & $M = \texttt{0x0000FFFF}$; \\
    & Внутренняя дивергенция: &
        $M_{\mathrm{tk}} = \texttt{0x000000FF}$,
        $M_{\mathrm{nt}} = \texttt{0x0000FF00}$ \\
    & & push $\{PC_C, \texttt{0x0000FFFF}, T\}$,
        push $\{PC_B, \texttt{0x0000FF00}, F\}$ \\
    & & $S = [\{100,M_0,T\},\{PC_D,M_{\mathrm{not\_outer}},F\},$ \\
    & & $\quad\quad \{PC_C,\texttt{0x0000FFFF},T\},\{PC_B,\texttt{0x0000FF00},F\}]$ \\
    \midrule
    3 & Исполнение $\mathrm{bb\_A}$ ($M = \texttt{0x000000FF}$)
      & $S$ без изменений \\
    \midrule
    4 & PC = $PC_B$ (CheckStack)
      & Снимаем $\{PC_B, \texttt{0x0000FF00}, F\}$: \\
    & & $M \leftarrow \texttt{0x0000FF00}$,
        $S = [\{100,M_0,T\},\{PC_D,M_{\mathrm{not}},F\},\{PC_C,\texttt{0x0000FFFF},T\}]$ \\
    \midrule
    5 & Исполнение $\mathrm{bb\_B}$ ($M = \texttt{0x0000FF00}$)
      & $S$ без изменений \\
    \midrule
    6 & PC = $PC_C$ (CheckStack)
      & Снимаем $\{PC_C, \texttt{0x0000FFFF}, T\}$: \\
    & & $M \leftarrow \texttt{0x0000FFFF}$,
        $S = [\{100, M_0, T\},\; \{PC_D, M_{\mathrm{not}}, F\}]$ \\
    \midrule
    7 & Исполнение $\mathrm{bb\_C}$ ($M = \texttt{0x0000FFFF}$)
      & $S$ без изменений \\
    \midrule
    8 & PC = $PC_D$ (CheckStack)
      & Снимаем $\{PC_D, M_{\mathrm{not}}, F\}$: \\
    & & $M \leftarrow \texttt{0xFFFF0000}$,
        $S = [\{100, M_0, T\}]$ \\
    \midrule
    9 & Исполнение $\mathrm{bb\_D}$ ($M = \texttt{0xFFFF0000}$)
      & $S$ без изменений \\
    \midrule
    10 & PC = 100 (CheckStack)
       & Снимаем $\{100, M_0, T\}$: \\
    & & $M \leftarrow \texttt{0xFFFFFFFF}$, $S = [\,]$ \\
    \midrule
    11 & Исполнение $\mathrm{bb\_E}$ ($M = \texttt{0xFFFFFFFF}$)
       & Полная реконвергенция \\
    \bottomrule
  \end{tabular}
\end{table}

Трассировка подтверждает, что LIFO-порядок снятия фреймов восстанавливает
маски в порядке «внутреннее-первым», затем «внешнее-первым», что
обеспечивает корректное исполнение всех ветвей при произвольной
глубине вложенности.
